%===================================================================
% Chapter 1: Introduction
% Phase-1 stub seeded from the thesis proposal. All prose here is a
% starting point -- revise freely. TODO markers indicate places that
% need your attention.
%===================================================================

%===================================================================
% Chapter 1: Introduction
%===================================================================
\chapter{Introduction}
\label{chap:intro}

Fifty years after its invention, SQL is more central to computing than it
has ever been. What began as a language for asking questions of a database
has become the implementation language of the modern data stack: the
dashboards an executive checks each morning, the feature tables that feed
machine learning models, the reports a regulator receives, and the metrics
behind customer-facing data products are all, underneath, SQL programs
executed against a data warehouse.

Two shifts distinguish today's SQL from the SQL of a decade ago. The first
is who, or rather what, writes it. A growing share of SQL is produced by
machines rather than people: business intelligence tools such as
Looker~\cite{lookergoogle} compile point-and-click analyses into queries
that routinely span hundreds of lines, a recent industry survey reports
that most analytics professionals now use AI assistance for code
development~\cite{dbt2025state}, and studies of BI-generated workloads
find scalar expressions and operator trees of a complexity no human would
write~\cite{getreal2018}. The second shift is how SQL is organized. Rather
than living as isolated statements, SQL increasingly accumulates into
recurring pipelines: version-controlled graphs of interdependent
transformations, re-executed on a schedule, in which each query consumes
the outputs of others. dbt, the most visible tool of this movement, is
used by tens of thousands of data teams every
week~\cite{dbtComplexity}, and the same dependency-structured pattern
appears in SQLMesh~\cite{sqlmeshOverview}, Dagster~\cite{dagsterAssets},
Airflow~\cite{airflowDags}, and the lakehouse architectures of every major
cloud vendor~\cite{databricksMedallion}.


These shifts unfold against a third, longer-running trend: the data
itself keeps growing. IDC estimates that the world created, captured,
and consumed more than 180 zettabytes of data in 2025, and projects
nearly 400 zettabytes by 2028, more than doubling in three
years~\cite{idcDataSphere2025}; an ever larger share of this data is
centralized in cloud platforms precisely so that it can be
queried~\cite{reinsel2018digitization}. Together, these trends have
turned the quality of SQL into an economic quantity. Cloud data
warehouses such as Snowflake~\cite{dageville2016snowflake} and
BigQuery~\cite{fernandes2015bigquery} bill computation by usage, so
the cost of a query grows with the data it scans, and an inefficient
query no longer merely runs slowly: over gigabytes it is an annoyance,
over terabytes it is a budget line, and inside a recurring pipeline it
is a budget line paid again on every refresh. Practitioners
consistently identify poorly written queries as a leading cause of
performance problems in
production~\cite{slow-query-impact-1,slow-query-impact-2,slow-query-impact-5}.
The trends above only sharpen the issue, because the SQL that machines
generate is precisely the SQL most likely to be written poorly. The
result is a widening gap between the complexity of the queries
organizations run, the volume of data they run them over, and the
quality of the queries they could be running.


% Together, these shifts have turned the quality of SQL into an economic
% quantity. Cloud data warehouses such as
% Snowflake~\cite{dageville2016snowflake} and
% BigQuery~\cite{fernandes2015bigquery} bill computation by usage, so an
% inefficient query no longer merely runs slowly: it costs money, and when
% that query sits inside a recurring pipeline, it costs money again on every
% refresh. Practitioners consistently identify poorly written queries as a
% leading cause of performance problems in
% production~\cite{slow-query-impact-1,slow-query-impact-2,slow-query-impact-5}.
% The trends above only sharpen the issue, because the SQL that machines
% generate is precisely the SQL most likely to be written poorly. The result
% is a widening gap between the complexity of the queries organizations run
% and the quality of the queries they could be running.

One might hope that query optimizers close this gap, since choosing
efficient execution strategies is exactly what they are for. But an optimizer operates within a boundary that is easy to overlook:
it selects the best execution plan for the query it is given, not the
best query its author could have written. Two formulations of the same
query, though semantically equivalent, can lead the same optimizer to
plans whose running times differ by orders of magnitude. When the
formulation is poor, the remedy must change the formulation. That
remedy is query rewriting: transforming a query into a semantically
equivalent form that the optimizer can turn into a faster plan.

The question is who performs this rewriting. The traditional answer is a
human expert, and for an important query in the hands of a seasoned
engineer this remains effective. But it does not scale. Rewriting a
single machine-generated query can require understanding hundreds of
lines of SQL together with the schema and constraints behind them, and
organizations accumulate such queries by the thousands. The automated
answer, developed over four decades of research and practice, is the
rewrite rule: a pattern that matches a fragment of a query and replaces
it with an equivalent, presumably faster counterpart. Rule-based
rewriting powers systems from Apache Calcite~\cite{begoli2018apache} to
recent tools that discover rules automatically~\cite{wetune}, and where
a rule applies, it works well. Its limitation is structural rather than
incidental: a rule-based rewriter is exactly as capable as its rule set,
and every rule must be anticipated, specified, and proven by someone in
advance. Queries whose patterns no rule anticipates receive no help at
all, and the machine-generated SQL now flooding data warehouses is an
endless source of unanticipated patterns.

Yet the rewrites themselves exist. Around any given query lies a vast
space of semantically equivalent formulations, and rules chart only the
few routes into that space that someone has already mapped. What has
been missing is a way to explore the space itself. This dissertation
argues that two developments make such exploration possible. Program
synthesis offers systematic search: given a way to score candidates and
check equivalence, it can traverse the space of SQL directly,
discovering rewrites no one has captured as a rule. Large language
models offer learned knowledge: trained on decades of accumulated code
and discussion, they have absorbed how experienced engineers reshape
queries, and can propose rewrites for exactly the complex, unfamiliar
queries where enumeration falters.

This dissertation develops techniques that use program synthesis and
large language models to rewrite SQL beyond predefined rules,
automatically discovering semantically equivalent and substantially
faster alternatives for individual queries and, ultimately, for entire
SQL pipelines.

\section{Rewriting as Search}
\label{sec:intro:search}

Rewriting without rules begins with a change of perspective. Instead of
asking which known transformations apply to a query, one can ask which
of the query's equivalent formulations runs fastest, and search for it
directly. Framed this way, query rewriting becomes a program synthesis
problem, and in principle any better formulation is discoverable,
whether or not anyone has thought to capture it as a rule. Two obstacles
stand between this idea and a working system. The space of equivalent
queries is unbounded, so a search that treats all candidates alike
drowns before it finds anything useful. And every candidate the search
does surface must be proven equivalent to the original, a problem that
is undecidable in general and expensive even in restricted forms.

SlabCity turns this idea into the first synthesis-based technique
capable of whole-query optimization without any rewrite rules. Its
central concept is the query dataflow, a semantic abstraction of how
data moves through a query. Dataflows give the search direction: they
induce a scoring function that organizes the space so that candidates
likely to be both equivalent and efficient are enumerated early, letting
the search restructure a query as a whole rather than patching isolated
fragments. Candidates that survive are admitted through layered
equivalence checking, in which inexpensive testing filters out the many
incorrect candidates and formal verification confirms the few that
remain, a spectrum of guarantees that Chapter~2 examines in detail.

The lesson of SlabCity is that rules are not necessary. Rule-free
rewriting is practical, and it discovers optimizations that rule-based
systems structurally cannot reach, including rewrites that no
composition of existing rules can express. But search has a ceiling of
its own, and it becomes visible as queries grow complex. The difficulty
is twofold. A complex query decomposes into many constituent dataflows,
and the number of ways to select and combine them into candidate
rewrites grows combinatorially, exploding the search space. At the same
time, the scoring function that steers the search is structural and
therefore coarse-grained: among this enormous number of candidates,
vast numbers receive similar scores, and the function can no longer
separate the promising few from the rest. What the search lacks at this
scale is not more candidates but semantic judgment about which rewrites
are worth attempting, the kind of judgment human experts accumulate
through experience.


\section{Rewriting with Learned Knowledge}
\label{sec:intro:knowledge}

Experienced engineers do not optimize a complex query by enumerating its
equivalent forms. They recognize it: this subquery recomputes what a
window function could compute once, that chain of self-joins is really a
single pass in disguise. Large language models, trained on decades of
public code and technical discussion, have absorbed a great deal of this
collective experience, and they offer what search lacks, namely
knowledge of how queries tend to be reshaped by people who know what
they are doing. What they do not offer is trust. A language model will confidently
propose rewrites that are subtly inequivalent, and rewrites that are
equivalent but no faster; and because commercial models charge for
every call, undisciplined use can spend more on the LLMs than the
resulting speedups save.

GenRewrite is the first holistic system to make language models
dependable for query rewriting, and its central idea is to let rules
return in a new form. Whenever a rewrite succeeds, GenRewrite distills
the insight behind it into a Natural Language Rewrite Rule (NLR2), a
human-readable strategy recorded together with the performance
bottleneck it resolves, and adds it to a growing repository. When a new
query arrives, GenRewrite inspects the query and its execution plan to
identify the likely bottleneck, retrieves the NLR2s that resolved
similar bottlenecks in earlier queries, and supplies them in the
model's prompt as concrete strategies to attempt. The same rules
explain each suggested rewrite in terms a human reviewer can check, and
they carry optimization knowledge from one query to the next, so the
system becomes more capable with every query it optimizes. When a
suggested rewrite turns out to be incorrect, a counterexample-guided
correction loop repairs its syntactic and semantic errors rather than
discarding it, salvaging rewrites that naive prompting would lose.

The rules that emerge are unlike the rules that came before. No one
wrote them in advance, they are stated in language rather than in
algebra, and they are applied with judgment rather than by pattern
matching. With them, GenRewrite optimizes the complex analytic queries
where both rule-based systems and synthesis-based search falter, and it
surpasses direct prompting while spending fewer model calls. What it
still shares with every rewriter before it is the unit of work: one
query at a time.

\section{From Queries to Pipelines}
\label{sec:intro:pipelines}

Modern analytics does not run one query at a time. It runs as
pipelines, in which hundreds of transformations feed one another on a
schedule, and this changes what optimization means. A query that is
perfectly written in isolation may still compute results that no
downstream consumer uses, duplicate work that a sibling transformation
already performs, or refresh hourly when its inputs change daily. No
single-query rewriter can see any of this, because the opportunities do
not live inside any one query. They live in the dependencies between
queries.

DAGSmith is built on the observation that those dependencies are not
merely execution constraints but optimization signals. Reading the
pipeline graph, it locates the places where dependency structure
indicates waste, uses a language model to propose refactorings that may
span multiple transformations, and admits a proposal only after
data-backed checking confirms that the pipeline's final outputs are
unchanged. Because a pipeline rewrite's value depends on how results
are stored and how often they are refreshed, DAGSmith weighs each
accepted rewrite jointly with materialization and refresh decisions,
using a learned cost model to select the combination that serves the
pipeline as a whole. The result is a pipeline that runs on the same
engine and the same orchestrator, produces the same outputs, and does
substantially less work to produce them.

DAGSmith completes the trajectory of this dissertation. Rewriting that
began with a single query and no rules ends at the scope where analytics
actually runs, with gains that no per-query technique, however strong,
can deliver alone.

\section{Contributions and Organization}
\label{sec:intro:contributions}

In summary, this dissertation makes the following contributions.

Chapter~3 presents SlabCity, the first synthesis-based technique for
whole-query SQL optimization that requires no rewrite rules. It
introduces query dataflows and a dataflow-driven scoring function that
make rule-free search practical, together with layered equivalence
checking that combines testing with bounded and full formal
verification. Evaluated on four workloads, including a newly curated
benchmark of more than 1,000 real-life queries, SlabCity optimizes more
queries than state-of-the-art rule-based rewriters, discovers rewrites
that no composition of their rules can express, and produces rewrites
that are often several times faster than the original queries.

Chapter~4 presents GenRewrite, the first holistic system for query
rewriting with large language models. It introduces Natural Language
Rewrite Rules as a mechanism for guiding, explaining, and transferring
rewriting knowledge, and a counterexample-guided correction loop that
repairs inequivalent candidates instead of discarding them. On TPC-DS,
GenRewrite improves 25 queries by a factor of two or more, 1.35 times
more than LLM-driven baselines and 2.6 times more than LLM-enhanced
rule-based systems, with consistent results on the Join Order Benchmark
and on machine-generated query variants.

Chapter~5 presents DAGSmith, which extends automatic rewriting from
individual queries to dependency-structured SQL pipelines. It
identifies pipeline dependencies as optimization signals, combines
LLM-guided refactoring with data-backed equivalence checking, and
jointly optimizes rewriting, materialization, and refresh frequency
under a learned cost model. On a real-world open-source healthcare
analytics pipeline of 255 transformations, DAGSmith reduces end-to-end
elapsed time by 42.6\% and warehouse compute cost by 67.7\%, far
exceeding both materialization-only tuning and state-of-the-art
single-query rewriting.

The remainder of this dissertation is organized as follows.
Chapter~2 provides background on the query rewriting problem, on
rule-based rewriting and its limits, on the spectrum of techniques for
establishing query equivalence, and on program synthesis and large
language models as they are used in this work. Chapters~3 through~5
present SlabCity, GenRewrite, and DAGSmith. Chapter~6 discusses
directions for future work, and Chapter~7 concludes.














\begin{comment}

\chapter{Introduction}
\label{chap:intro}

Over the past decade, data volumes and data infrastructure have both
exploded. Modern organizations increasingly express their analytics,
reporting, and data transformation logic in SQL, whether written by
hand, generated by business-intelligence tools, or assembled into
large pipelines of interdependent models. Yet much of this SQL is far
from optimal: poorly written queries are a major cause of slow
performance, and in cloud databases they translate directly into
excessive cost. Traditional query optimizers rely on carefully crafted
rewrite rules, but maintaining and extending these rules as systems,
hardware, and workloads change requires substantial expert effort, and
many performance issues still end up being diagnosed and rewritten
manually by senior engineers or DBAs, one problematic query or
pipeline at a time.

This dissertation asks whether we can push SQL optimization beyond
existing rules, using program synthesis and large language models
(LLMs) to automatically discover and apply rewrites that human experts
would otherwise have to craft by hand. The goal is not to replace
traditional optimizers, but to surround them with higher-level
rewriters that (i)~search over a richer space of semantically
equivalent formulations, (ii)~reuse knowledge about successful
rewrites across queries and workloads, and (iii)~scale from individual
queries to large, dependency-rich pipelines.

% TODO: consider adding an explicit, one-sentence thesis statement
% here (see, e.g., how other UM dissertations phrase "Thesis
% Statement" as its own subsection).

\section{Overview of the Dissertation}
\label{sec:intro:overview}

To explore this space, we develop three systems that form a coherent
trajectory from synthesis to LLMs, and from single queries to full
pipelines.

\paragraph{SlabCity (\cref{slab:chap}).}
SlabCity is a whole-query, synthesis-based SQL rewriter that searches
for semantically equivalent, faster queries without relying on
predefined rules. It defines \emph{dataflows} for SQL queries and uses
them to stratify the synthesis search space via a novel dataflow-based
scoring function that focuses the search on promising rewrites.
SlabCity employs a counterexample-guided inductive synthesis (CEGIS)
loop that starts with lightweight tests and bounded checks, resorting
to SMT verification only when needed. On curated LeetCode and Calcite
workloads, it rewrites more queries and achieves multi-fold latency
reductions compared to rule-based systems.

\paragraph{GenRewrite (\cref{gen:chap}).}
GenRewrite targets the same single-query setting but replaces
synthesis-based search with LLM-guided generation anchored by
Natural-Language Rewrite Rules (NLR2s). NLR2s capture successful
rewrite strategies in human-readable form and serve both as prompts to
the LLM and as a medium for transferring knowledge across queries:
once distilled from an effective rewrite, an NLR2 can be reused to
improve relevant future queries, allowing GenRewrite to become more
capable over time. For each query, GenRewrite consults its growing
NLR2 repository, uses plan analysis to retrieve bottleneck-relevant
rules, and feeds them to the LLM, which proposes and iteratively
refines candidate rewrites via a semantic-then-syntax correction loop.
Refined candidates are then passed to off-the-shelf testers and
verifiers to filter out inequivalent rewrites. Across TPC-DS and JOB
(including their SQLStorm-generated variants), GenRewrite attains
higher coverage and larger speedups than both rule-based systems and
baseline LLM prompting.

\paragraph{DAGSmith (\cref{dag:chap}).}
DAGSmith lifts these ideas from single queries to pipelines of
interdependent SQL models, as popularized by dbt-style data
transformation frameworks. Given a workload represented as a query
DAG, DAGSmith analyzes compiled SQL and dependency edges to identify
regions likely to contain optimization opportunities, leverages LLMs
to reason over the relevant subgraph and propose concrete
refactorings, and separates proposal, criticism, SQL generation, and
data-backed equivalence checking so unsafe rewrites can be rejected
before adoption. Adopted candidates are co-optimized with
materialization decisions using a learned cost model and an iterative
ILP formulation, and frequency information reweights cost estimates to
detect over-scheduled work. The result is a rewritten,
output-equivalent SQL pipeline that remains executable by the same SQL
engine and orchestration framework, but is significantly more
performant.

% TODO: the proposal referred to this system as DAGWeaver; the final
% system is DAGSmith. Double-check that no stray "DAGWeaver" mentions
% remain anywhere in the final dissertation text.

\section{Contributions}
\label{sec:intro:contributions}

% TODO: Phase-2 material. Summarize the combined contributions of the
% three systems here, or fold this into the overview above. The
% per-paper contribution lists currently live in each chapter's
% introduction section.
In summary, this dissertation traces a path from synthesis to LLMs for
SQL query rewriting beyond rules, delivering substantial performance
gains from individual queries to large, dependency-rich SQL pipelines.

\section{Organization}
\label{sec:intro:organization}

The remainder of this dissertation is organized as follows.
\Cref{slab:chap} presents SlabCity, our synthesis-based whole-query
rewriter. \Cref{gen:chap} presents GenRewrite, which pivots to
LLM-guided rewriting with natural-language rewrite rules.
\Cref{dag:chap} presents DAGSmith, which extends rewriting to full SQL
pipeline DAGs. \Cref{chap:future} discusses directions for future
work, and \cref{chap:conclusion} concludes.
\end{comment}